\section{Related Work}
\label{sec:related}

\paragraph{PIM and heterogeneous acceleration for LLM inference.}
Recent systems work argues that LLM decoding contains kernels with heterogeneous bottlenecks and that PIM or near-memory execution can improve performance in the right regimes. PAPI is the closest direct prior in spirit: it studies dynamic kernel assignment for LLM decoding on a PIM-enabled heterogeneous system and shows that static mapping is insufficient in the presence of dynamic parallelism \citep{he2025papi}. Other architecture papers such as PIM Is All You Need and CENT push the argument further by proposing concrete PIM-enabled or GPU-free systems for LLM inference \citep{gu2025pim,wang2025cent}. These papers are important because they motivate why PIM may matter at decode time. Our work differs in scope and claim boundary. We do not propose a new PIM architecture or claim real hardware speedups. Instead, we use runtime-grounded traces and a virtual backend to characterize \emph{when} PIM-style offloading appears useful and which runtime signals best predict regime entry under that backend.

\paragraph{KV-cache-centric inference systems.}
Another closely related line of work treats the KV cache as the central systems bottleneck in long-context inference. Mooncake, for example, proposes a KVCache-centric disaggregated serving architecture and scheduler for long-context LLM service \citep{qin2025mooncake}. Other work studies dynamic KV-cache placement or online scheduling under heterogeneous memory constraints \citep{akhauri2025kvplacement,microsoft2025kvscheduling}. These papers substantially overlap with our mechanism story: KV traffic and memory pressure matter. However, they ask a different question. Their main objective is serving architecture, memory placement, or online scheduling under explicit storage constraints, whereas our focus is a trace-driven characterization of virtual PIM offload regimes during decoding. In particular, we ask whether KV-related signals provide better predictive support for positive-regime detection than a static attention rule.

\paragraph{LLM inference scheduling and serving optimization.}
Constraint-aware and throughput-aware LLM schedulers also provide important context. ExeGPT studies resource scheduling for LLM inference under latency constraints and shows that execution configuration should adapt to workload characteristics \citep{oh2024exegpt}. These works support the broader view that LLM inference should be treated as a scheduling problem rather than a fixed execution template. Our policy analysis is aligned with that perspective, but we focus on a narrower issue: not the entire serving system, but the emergence of positive virtual PIM offload regimes at decode time. This is why our analysis emphasizes regime entry, negative short-context cases, and runtime signals such as memory pressure proxy rather than only throughput-oriented scheduler optimization.

\paragraph{Trace-driven and simulation-based LLM evaluation.}
A growing body of work uses simulation or replay to study LLM inference systems without rebuilding full production stacks. This makes trace-driven methodology itself a crowded space, which means our novelty does not come from using simulation alone. Our contribution is instead the combination of runtime-grounded decode traces, a PIM-specific virtual what-if backend, and an empirical finding: the onset of virtual offload value is context-sensitive, model-family dependent, and only partially captured by static attention labels. We therefore position the paper as an empirical characterization and decision-signal study enabled by a trace-driven framework, rather than as a general LLM simulator paper.
